% !TEX root = ../main.tex

\chapter{实验结果与讨论}
\chaptereng{Experimental Results and Discussion}


\section{平面炉初步温度测量}

在完成试验台搭建和激光器波长锁定后，首先采用甲烷/空气平面火焰进行初步的TDLAS测量，以验证系统在真实火焰条件下的工作性能并确定最佳测量工况。

\subsection{实验准备与工况参数}

实验前完成以下准备工作：完成实验方案设计、工况参数确定和安全预案制定。采用甲烷空气火焰（8 SLM CH$_4$ / 50 SLM air，当量比$\phi=1.0$），激光器参数设置为50.9 mA、$16^\circ$C、200 Hz、0.8 Vpp，该组参数基于第三章光谱仿真确定的波长位置进行微调后获得。

\subsection{吸收峰验证}

图~\ref{fig:absorption_peaks} 展示了甲烷/空气平面炉火焰的典型吸收光谱。

\begin{figure}[!htp]
	\centering
	\includegraphics[width=0.7\textwidth]{image37.png}
	\bicaption{平面炉火焰中H$_2$O吸收光谱}{H$_2$O absorption spectrum in flat-flame burner flame}
	\label{fig:absorption_peaks}
\end{figure}

测试结果成功观测到位于7153.75 cm$^{-1}$和7154.35 cm$^{-1}$的两个水吸收峰，与第三章Flame Spectrum仿真结果吻合良好。该结果验证了以下三点：第一，激光器波长锁定准确，在200 Hz扫描频率下可稳定覆盖目标双线；第二，光谱信噪比满足后续Voigt拟合对原始数据质量的要求；第三，甲烷/空气平面炉能够产生稳定的火焰环境，适合作为TDLAS测量的火源。上述吸收峰的明确识别为后续高频扫描实验确定了激光器工作参数。


\section{平面炉高频扫描实验}

在初步验证吸收峰位置后，需进一步提升扫描频率以适应实际火焰测量需求。

\subsection{高频扫描的必要性}

火焰的内部温度、组分浓度和压力在时间和空间上都存在剧烈的脉动。若扫描频率低，单次扫描耗时长，在这期间火焰状态早已改变，得到的吸收光谱实际上是一段时间内不同火焰状态的模糊平均，由此反演出的温度和浓度就会严重失真，尤其会让温度测量值偏低、浓度分布展宽。

TDLAS通过拟合吸收谱线的线型来反演温度和浓度等参数，这要求单次扫描内测量条件基本稳定。在低频扫描时，激光波长缓慢扫过吸收线，而这期间火焰的温度/压力变化会导致谱线中心频率漂移和展宽抖动，最终测得的谱线是这些动态过程的叠加，线型发生畸变，反演误差增大。

因此，需要将扫描频率提升至kHz级别。本实验将扫描频率提升至\textbf{10 kHz}，测量火焰上方1--5 mm高度，成功采集了高频条件下的吸收光谱。

\subsection{实验设置与结果}

高频扫描实验采用与初步实验相同的火焰工况（8 SLM CH$_4$ / 50 SLM air，$\phi=1.0$），激光器参数调整至：50.9 mA、16.1$^\circ$C、10 kHz、0.4 Vpp。伺服平台控制炉体在X方向$-15$至$+15$ mm、Y方向1至5 mm范围内步进扫描。

实验成功采集了10 kHz扫描频率下的吸收光谱，验证了TDLAS系统在高频模式下的稳定工作能力。初步观测结果显示，高频扫描有效避免了因火焰脉动导致的谱线畸变问题，为后续精确温度反演奠定了数据基础。


\section{带伴流平面炉扫描实验}

尽管平面炉本身具有较好的火焰稳定性，但开放环境下火焰边界仍容易受环境气流和热浮力扰动的影响。

\subsection{伴流的作用机制}

在开放火焰中，火焰边界容易受气流、热浮力扰动影响而产生抖动或边缘卷吸。同轴伴流可以产生一个低速、平稳的气流保护层，有效抑制外界紊乱对火焰本体的干扰，并减少边缘的空气卷吸。

此外，开放火焰的边缘会卷吸空气中的氧气，这会导致在火焰边界形成一个当量比无法控制、温度和组分梯度大的反应层。用惰性气体伴流（通常为N$_2$），可阻断氧气扩散进入火焰高温区，消除火焰边缘的异常，使火焰产物的组分浓度和温度严格由实验设定燃料燃烧决定，而非受边界二次反应干扰。这对于高精度的火焰标定至关重要。

\subsection{实验结果}

采用带N$_2$伴流（伴流流量20 SLM）的平面炉进行甲烷/空气火焰扫描实验，在与高频扫描实验相同工况下（8 SLM CH$_4$ / 50 SLM air，$\phi=1.0$，10 kHz扫描频率）采集吸收光谱。初步结果表明，伴流的加入显著改善了火焰边缘的温度均匀性，边缘区域的温度梯度较无伴流工况有所减小。